\section{Instalación de Docker}

\subsection{Inspección y adaptación a Fedora}

Las instrucciones de referencia presentan procedimientos para Debian y
Ubuntu. En este equipo no resultaba correcto repetirlos, pues Docker ya se
encontraba instalado mediante los paquetes propios de Fedora:
\texttt{moby-engine}, \texttt{docker-cli}, \texttt{containerd} y
\texttt{docker-buildx}. Por tanto, el procedimiento se concentró en comprobar
los paquetes, habilitar el servicio y validar el funcionamiento del motor.

El servicio se configuró para iniciarse automáticamente y se arrancó en la
sesión actual:

\begin{lstlisting}[style=terminal]
$ sudo systemctl enable --now docker
$ systemctl is-enabled docker
enabled
$ systemctl is-active docker
active
\end{lstlisting}

El usuario ya figuraba como integrante del grupo \texttt{docker}; sin embargo,
la sesión desde la que se inició la práctica todavía no había incorporado esa
pertenencia. Para aplicar el grupo sin reiniciar el equipo se utilizó una
subsesión equivalente a volver a iniciar sesión:

\begin{lstlisting}[style=terminal]
$ getent group docker
docker:x:971:omar
$ sg docker -c "docker version"
\end{lstlisting}

\begin{nota}{Implicación del grupo docker}
La pertenencia al grupo \texttt{docker} permite controlar el demonio y, en la
práctica, concede capacidades comparables con privilegios administrativos. El
acceso debe asignarse únicamente a usuarios de confianza.
\end{nota}

\subsection{Prueba funcional}

La comunicación entre el cliente y el servidor se comprobó con sus versiones.
Después se ejecutó la imagen mínima \texttt{hello-world}. Docker descargó la
imagen, creó un contenedor efímero y devolvió el mensaje de confirmación, lo que
demostró el funcionamiento completo de la ruta cliente--demonio--registro.

\begin{lstlisting}[style=terminal]
$ docker version
$ docker run --rm hello-world
\end{lstlisting}

\evidencia{02_evidencia.png}{Validación del cliente, el demonio y la ejecución de la imagen \texttt{hello-world}.}
